\documentclass[a4paper, 10pt]{article}
\usepackage[utf8]{inputenc}
\usepackage[T1]{fontenc}
\usepackage[french]{babel}
\usepackage{graphicx}
\usepackage[a4paper,margin=2.45cm]{geometry}
\usepackage{amsmath}
\numberwithin{equation}{section}


\renewcommand{\thesection}{\Roman{section}}

\title{ \bf Supersymmetric and Shape-Invariant Generalization for
Non-resonant Jaynes-Cummings Systems}
\author{A. N. F. Aleixo \footnote{Electronic address: aleixo@nucth.physics.wisc.edu} \\
\it Department of Physics, University of Wisconsin \\
\it Madison, Wisconsin 53706 USA  \footnote{Permanent address: Instituto de F´ısica, Universidade Federal do Rio de Janeiro, RJ - Brazil.}
\and \\
A. B. Balantekin \footnote{Electronic address: baha@nucth.physics.wisc.edu}\smallskip, \\
\it Max-Planck-Institut für Kernphysik, Postfach 103980, D-69029 Heidelberg, Germany \footnote{Permanent address:Department of Physics, University of Wisconsin, Madison, Wisconsin 53706
USA}
\and \\ 
M. A. Cândido Ribeiro \footnote{Electronic address: macr@df.ibilce.unesp.br} \smallskip, \\
\it Departamento de F\'isica - Instituto de Biociências, Letras e Ciências Exatas \\
UNESP, S\~ao José do Rio Preto, SP - Brazil}
\date{(November 13, 2018)}

\begin{document}

\maketitle
\reversemarginpar
\marginpar{\rotatebox{90}{\Large arxiv.org/abs/quant-ph/0005045v1}}

\begin{center}
    \section*{Abstract}
\end{center}

A class of shape-invariant bound-state problems which represent transitions
in a two-level system introduced earlier are generalized to include arbitrary
energy splittings between the two levels. We show that the coupled-channel
Hamiltonians obtained correspond to the generalization of the non-resonant
Jaynes-Cummings Hamiltonian, widely used in quantized theories of laser.
In this general context, we determine the eigenstates, eigenvalues, the time
evolution matrix, and the population inversion matrix factor.\\

\begin{flushright}
    Typeset using \LaTeX
\end{flushright}

\newpage
\begin{center}
     \bf \section{INTRODUCTION}
\end{center}

The integrability condition called shape-invariance originates in supersymmetric quantum mechanics ~\cite{[1,2]}. The separable positive-definite Hamiltonian $ \hat{H_1} = \hat{A}^{\dagger} \hat{A} $ is called shape-invariant if the condition 
\begin{equation}
     \hat{A}(a_1) \hat{A}^{\dagger}(a_1) = \hat{A}^{\dagger}(a_2) \hat{A}(a_2) + R(a_2) 
\end{equation}
is satisfied ~\cite{[3]}. In this equation $a_1$ and $a_2$ represent parameters of the Hamiltionian. The parameter $a_2$ is a function of $a_1$ and the remainder $R(a_1)$ is independent of the dynamical variables such as position and momentum. Even though not all exactly-solvable problems are shape-invariant ~\cite{[4]}, shape invariance, especially in its algebraic formulation ~\cite{[5-7]}, has prove, to be a powerful technique to study exactly-solvable systems.\\
\vspace{0,5cm}

\begin{center}
    {\bf\section{THE GENERALIZED NON-RESONANT JAYNES-CUMMINGS HAMILTONIAN}}
\end{center}

\subsection*{The Resonant Limit}

\subsection*{The Standard Jaynes-Cummings Limit}
\vspace{0,5 cm}

\begin{center}
    {\bf\section{THE TIME EVOLUTION OF THE SYSTEM}}
\end{center}

\noindent
To study the time-dependent Schrödinger equation for a Jaynes-Cummings system in
non-resonant situation
\begin{equation}
    i\hbar\frac{\partial}{\partial t}|\Psi(t)\rangle = ({\bf \hat{H}}_o + {\bf \hat{H}}_{int}) |\Psi(t)\rangle
\end{equation}
\noindent
we can write the wavefunction as
\begin{equation}
    |\Psi(t)\rangle = exp(-i{\bf \hat{H}}_ot/\hbar) |\Psi_i(t)\rangle ,
\end{equation}

\noindent
and, by substituting this into Schrödinger equation and taking into account the commutation property between ${\bf\hat{H}}_o$ et ${\bf \hat{H}}_{int}$ , we obtain

\begin{equation}
    i\hbar\frac{\partial}{\partial t}|\Psi_i(t)\rangle = {\bf \hat{H}}_{int}|\Psi_i(t)\rangle
\end{equation}
\noindent
We introduce the evolution matrix ${\bf \hat{U}}_i(t,0)$ :
\begin{equation}
    |\Psi_i(t)\rangle = {\bf \hat{U}}_i(t,0)|\Psi_i(0)\rangle
\end{equation}

\noindent
which satisfies the equation
\begin{equation}
    i\hbar\frac{\partial}{\partial t} {\bf \hat{U}}_i(t,0) = {\bf \hat{H}}_{int} {\bf \hat{U}}_i(t,0) ,
\end{equation}

\noindent
that is, in matrix form, written as
\begin{equation}
i\hbar \left [\begin{array}{cc}
      \hat{U}'_{11}   & \hat{U}'_{12}  \\
      \hat{U}'_{21}& \hat{U}'_{22}
    \end{array} \right] = \alpha \left[\begin{array}{cc}
     \beta    & \hat{T}\hat{B}_-  \\
      \hat{B}_+\hat{T}^\dagger   & -\beta 
    \end{array}\right] \left [\begin{array}{cc}
      \hat{U}_{11}   & \hat{U}_{12}  \\
    \hat{U}_{21} & \hat{U}_{22}
    \end{array} \right]
\end{equation} 

\noindent
where the primes denote the time derivative. One fast way to diagonalize the evolution
matrix differential equation is by differentiating Eq. (3.5) with respect to time. We find

\begin{equation}
    i\hbar\frac{\partial^2}{\partial t^2}{\bf \hat{U}}_i(t,0) = {\bf \hat{H}}_{int}\frac{\partial}{\partial t}{\bf\hat{U}}_i(t,0) = \frac{1}{i\hbar}{\bf \hat{H}}_{int}^2(t,0)
\end{equation}
\noindent
which can be written as
\begin{equation}
    \left [\begin{array}{cc}
   \hat{U''}_{11}  & \hat{U''}_{12} \\
    \hat{U''}_{21} & \hat{U''}_{22}
\end{array} \right] = -\left[\begin{array}{cc}
    \omega_1 & 0  \\
    0 & \omega_2
\end{array} \right] \left[\begin{array}{cc}
   \hat{U}_{11}  &  \hat{U}_{12} \\
    \hat{U}_{21} & \hat{U}_{22}
\end{array} \right]
\end{equation}

\noindent
where

\begin{equation}
    \hbar\hat{\omega}_1 = \alpha \sqrt{\hat{T}\hat{B}_-\hat{B}_+ +\beta^2} = \sqrt{\hbar\Omega \hat{H}_2 + (\hbar\Delta)^2} 
\end{equation}

\begin{equation}
    \hbar\hat{\omega}_2 = \alpha \sqrt{\hat{B}_-\hat{B}_+ +\beta^2} = \sqrt{\hbar\Omega \hat{H}_1 + (\hbar\Delta)^2}
\end{equation}

\noindent
Since by initial conditions ${\bf \hat{U}}_i(0,0) = \hat{I}$, , then we can write the solution of the evolution matrix differential equation (3.7) as

\begin{equation}
    {\bf \hat{U}}_i(t,0) = \left [\begin{array}{cc}
      \cos(\omega_1 t)   & \sin(\omega_1 t)\hat{C} \\
        \sin(\omega_2 t)\hat{D} & \cos(\omega_2 t)
    \end{array} \right]
\end{equation}

\noindent
and the $\hat{C}$ the $\hat{D}$ operators can be determined by the unitarity conditions

\begin{equation}
    {\bf \hat{U}}_i^{\dagger}(t,0){\bf \hat{U}}_i(t,0) = {\bf \hat{U}}_i(t,0){\bf \hat{U}}_i^ {\dagger}(t,0) = {\bf \hat{I}} .
\end{equation}

\begin{thebibliography}{9}
\bibitem{[1]} E. Witten, Nucl. Phys. B 185, 513 (1981).
\bibitem{[2]} For a recent review see F. Cooper, A. Khare and U. Sukhatme, Phys. Rep. 251, 267
(1995).
\bibitem{[3]} L. Gendenshtein, Pis’ma Zh. Eksp. Teor. Fiz. 38, 299 (1983) [JETP Lett. 38, 356
(1983)].
\bibitem{[4]} F. Cooper, J. N. Ginocchio and A. Khare, Phys. Rev. D 36, 2458 (1987).
\bibitem{[5]} A. B. Balantekin, Phys. Rev. A 57, 4188 (1998).
\bibitem{[6]} S. Chaturvedi, R. Dutt, A. Gangopadhyaya, P. Panigrahi, C. Rasinariu, and U.
Sukhatme, Phys. Lett. A 248, 109 (1998).
\bibitem{[7]} A. N. F. Aleixo, A. B. Balantekin, and M. A. Cândido Ribeiro, J. Phys. A: Math. Gen.
32, 2785 (1999).
\bibitem{[8]} A. B. Balantekin, M. A. Cândido Ribeiro, and A. N. F. Aleixo, J. Phys. A: Math. Gen.
33, 1503 (2000).
\end{thebibliography}
\end{document}
